\documentclass{article}
\usepackage{amsmath, amssymb, amsthm}
\usepackage{geometry}
\geometry{a4paper, margin=1in}

\title{Discussion on \( W^* \)-Algebras and Convexity}
\author{}
\date{}

\begin{document}

\maketitle

\section*{Question 1: If \( f \) is a linear functional defined on a convex set, does its supremum occur at extreme points of this set?}

Yes, if \( f \) is a linear functional defined on a convex set \( C \), then the supremum (sup) of \( f \) over \( C \) occurs at an extreme point of \( C \), provided that \( C \) is compact (closed and bounded in finite dimensions) and \( f \) is continuous. This result is a consequence of the \textbf{Krein-Milman Theorem} and the properties of linear functionals on convex sets.

\subsection*{Explanation:}
\begin{enumerate}
    \item \textbf{Linear Functional}: A linear functional \( f \) is a linear map from a vector space to the field of scalars (e.g., \( \mathbb{R} \) or \( \mathbb{C} \)). For a convex set \( C \), \( f \) can be written as \( f(x) = \langle a, x \rangle \) for some vector \( a \), where \( \langle \cdot, \cdot \rangle \) is an inner product.

    \item \textbf{Extreme Points}: An extreme point of a convex set \( C \) is a point that cannot be expressed as a convex combination of other points in \( C \). For example, in a convex polygon, the extreme points are the vertices.

    \item \textbf{Supremum at Extreme Points}: Since \( f \) is linear, it attains its maximum (or supremum) on \( C \) at an extreme point of \( C \). This is because:
    \begin{itemize}
        \item The maximum of a linear functional over a convex set must occur on the boundary of the set.
        \item On a compact convex set, the boundary is generated by the extreme points (by the Krein-Milman Theorem).
    \end{itemize}

    \item \textbf{Compactness Assumption}: If \( C \) is not compact, the supremum of \( f \) might not be attained at all. However, if \( C \) is compact, the supremum is achieved, and it occurs at an extreme point.
\end{enumerate}

\subsection*{Example:}
Consider a convex polygon in \( \mathbb{R}^2 \). A linear functional \( f(x, y) = ax + by \) will attain its maximum value at one of the vertices (extreme points) of the polygon.

\subsection*{Conclusion:}
If \( f \) is a linear functional defined on a compact convex set \( C \), then the supremum of \( f \) over \( C \) is achieved at an extreme point of \( C \). This is a fundamental result in convex analysis and optimization.

\section*{Question 2: Is this true for infinite convex sets?}

The statement that the supremum of a linear functional \( f \) over a convex set \( C \) occurs at an extreme point of \( C \) is \textbf{not always true for infinite convex sets}, even if \( C \) is compact. The behavior depends on the specific properties of the set \( C \) and the linear functional \( f \). Let's break this down:

\subsection*{1. Finite-Dimensional Case (Compact Convex Sets):}
\begin{itemize}
    \item If \( C \) is a \textbf{compact convex set in a finite-dimensional space} (e.g., \( \mathbb{R}^n \)), then the Krein-Milman Theorem guarantees that \( C \) is the closed convex hull of its extreme points.
    \item For a linear functional \( f \), the supremum of \( f \) over \( C \) is achieved at an extreme point of \( C \). This is because \( f \) is continuous, and the maximum of a continuous function over a compact set is attained at an extreme point.
\end{itemize}

\subsection*{2. Infinite-Dimensional Case:}
\begin{itemize}
    \item In infinite-dimensional spaces (e.g., function spaces or spaces of sequences), the situation is more nuanced. Even if \( C \) is compact and convex, the Krein-Milman Theorem still holds (i.e., \( C \) is the closed convex hull of its extreme points), but the supremum of \( f \) over \( C \) may \textbf{not necessarily occur at an extreme point}.
    \item The key issue is that in infinite dimensions, compactness does not imply that the set is the convex hull of its extreme points in a "strong" sense. Extreme points may not be sufficient to characterize the behavior of linear functionals.
\end{itemize}

\subsection*{3. Counterexample in Infinite Dimensions:}
Consider the following example in the space \( \ell^2 \) (the space of square-summable sequences):
\begin{itemize}
    \item Let \( C \) be the closed unit ball in \( \ell^2 \), which is a compact convex set in the weak topology.
    \item Define a linear functional \( f \) on \( \ell^2 \) by \( f(x) = \sum_{n=1}^\infty \frac{1}{n} x_n \), where \( x = (x_1, x_2, \dots) \).
    \item The supremum of \( f \) over \( C \) is \( \sum_{n=1}^\infty \frac{1}{n^2} = \frac{\pi^2}{6} \), but this supremum is \textbf{not attained at any extreme point of \( C \)}. The extreme points of \( C \) are the unit vectors \( e_n = (0, \dots, 0, 1, 0, \dots) \), but \( f(e_n) = \frac{1}{n} \), which does not achieve the supremum.
\end{itemize}

\subsection*{4. When Does It Hold in Infinite Dimensions?}
\begin{itemize}
    \item If \( C \) is a \textbf{compact convex set in a locally convex topological vector space}, and \( f \) is a continuous linear functional, then the supremum of \( f \) over \( C \) is achieved at an extreme point of \( C \) \textbf{if and only if} the set of points where \( f \) attains its supremum is a face of \( C \) that contains an extreme point.
    \item In some special cases (e.g., when \( C \) is the closed convex hull of its extreme points and satisfies additional conditions), the supremum of \( f \) will occur at an extreme point.
\end{itemize}

\subsection*{Conclusion:}
For \textbf{infinite convex sets}, the supremum of a linear functional \( f \) does not necessarily occur at an extreme point, even if the set is compact. The result depends on the structure of the set and the space in which it is embedded. In finite dimensions, the statement holds under compactness, but in infinite dimensions, additional conditions are required.

\section*{Question 3: Is the set of positive elements of a \( W^* \)-algebra a convex set?}

Yes, the set of positive elements in a \( W^* \)-algebra (also known as a von Neumann algebra) is a \textbf{convex set}. Let's break this down:

\subsection*{1. Positive Elements in a \( W^* \)-Algebra:}
\begin{itemize}
    \item A \( W^* \)-algebra is a \( C^* \)-algebra that is also a dual space (i.e., it has a predual).
    \item An element \( x \) in a \( W^* \)-algebra \( \mathcal{A} \) is called \textbf{positive} if it is self-adjoint (\( x = x^* \)) and its spectrum \( \sigma(x) \) is contained in \( [0, \infty) \). We write \( x \geq 0 \) to denote that \( x \) is positive.
\end{itemize}

\subsection*{2. Convexity of the Set of Positive Elements:}
\begin{itemize}
    \item The set of positive elements in \( \mathcal{A} \), denoted \( \mathcal{A}_+ \), is convex. This means that for any two positive elements \( x, y \in \mathcal{A}_+ \) and any \( \lambda \in [0, 1] \), the convex combination \( \lambda x + (1 - \lambda) y \) is also positive.
    \item \textbf{Proof}:
    \begin{itemize}
        \item Since \( x \) and \( y \) are positive, \( x = x^* \), \( y = y^* \), and their spectra are non-negative.
        \item The convex combination \( \lambda x + (1 - \lambda) y \) is self-adjoint because:
        \[
        (\lambda x + (1 - \lambda) y)^* = \lambda x^* + (1 - \lambda) y^* = \lambda x + (1 - \lambda) y.
        \]
        \item The spectrum of \( \lambda x + (1 - \lambda) y \) is non-negative because the set of positive operators is closed under non-negative linear combinations.
    \end{itemize}
\end{itemize}

\subsection*{3. Additional Properties:}
\begin{itemize}
    \item The set \( \mathcal{A}_+ \) is also a \textbf{cone}, meaning that if \( x \in \mathcal{A}_+ \) and \( \alpha \geq 0 \), then \( \alpha x \in \mathcal{A}_+ \).
    \item \( \mathcal{A}_+ \) is closed in the weak-\( * \) topology (and hence in the norm topology) because the positive elements are characterized by conditions involving the weak-\( * \) continuous functional calculus.
\end{itemize}

\subsection*{4. Geometric Interpretation:}
\begin{itemize}
    \item The set \( \mathcal{A}_+ \) forms a "positive cone" in the \( W^* \)-algebra \( \mathcal{A} \). It is a convex cone because it is closed under convex combinations and non-negative scalar multiplication.
\end{itemize}

\subsection*{Conclusion:}
The set of positive elements in a \( W^* \)-algebra is indeed a \textbf{convex set}. It is also a convex cone, and it plays a fundamental role in the structure and analysis of \( W^* \)-algebras.

\section*{Question 4: Is a \( W^* \)-algebra a locally convex topological vector space?}

Yes, a \( W^* \)-algebra (also known as a von Neumann algebra) is a \textbf{locally convex topological vector space}. Let's explain why this is the case and clarify the relevant concepts:

\subsection*{1. What is a \( W^* \)-Algebra?}
\begin{itemize}
    \item A \( W^* \)-algebra is a \( C^* \)-algebra that is also a dual space. That is, it is the dual of some Banach space, called its \textbf{predual}.
    \item Equivalently, a \( W^* \)-algebra is a \( C^* \)-algebra that is closed in the weak-\( * \) topology induced by its predual.
\end{itemize}

\subsection*{2. Locally Convex Topological Vector Space:}
\begin{itemize}
    \item A \textbf{locally convex topological vector space} is a vector space equipped with a topology that is generated by a family of seminorms. This topology ensures that the space has a rich supply of continuous linear functionals and convex neighborhoods of the origin.
    \item Key properties of locally convex spaces include:
    \begin{itemize}
        \item The topology is Hausdorff.
        \item The topology is generated by convex open sets.
        \item The Hahn-Banach theorem holds, allowing for the separation of convex sets by continuous linear functionals.
    \end{itemize}
\end{itemize}

\subsection*{3. Why is a \( W^* \)-Algebra Locally Convex?}
\begin{itemize}
    \item A \( W^* \)-algebra \( \mathcal{A} \) is a Banach space (since it is a \( C^* \)-algebra) and is also a dual space (by definition). As a dual space, it carries the \textbf{weak-\( * \) topology}, which is a locally convex topology.
    \item The weak-\( * \) topology on \( \mathcal{A} \) is generated by the seminorms:
    \[
    p_\phi(x) = |\phi(x)|, \quad \text{for } x \in \mathcal{A},
    \]
    where \( \phi \) ranges over the predual of \( \mathcal{A} \). These seminorms define a locally convex topology.
    \item The weak-\( * \) topology is coarser than the norm topology but still makes \( \mathcal{A} \) a locally convex topological vector space.
\end{itemize}

\subsection*{4. Additional Properties:}
\begin{itemize}
    \item The weak-\( * \) topology on a \( W^* \)-algebra is Hausdorff, meaning that distinct points can be separated by open sets.
    \item The algebra operations (addition, scalar multiplication, and multiplication) are continuous in the weak-\( * \) topology, making \( \mathcal{A} \) a topological algebra.
    \item The positive cone \( \mathcal{A}_+ \) is closed in the weak-\( * \) topology, and the order structure of \( \mathcal{A} \) is compatible with this topology.
\end{itemize}

\subsection*{5. Comparison with Other Topologies:}
\begin{itemize}
    \item A \( W^* \)-algebra also carries the \textbf{norm topology}, which is finer than the weak-\( * \) topology. In the norm topology, \( \mathcal{A} \) is a Banach space and hence also a locally convex space.
    \item However, the weak-\( * \) topology is particularly important in the study of \( W^* \)-algebras because it reflects the duality with the predual and is used to define important concepts like normal states and normal representations.
\end{itemize}

\subsection*{Conclusion:}
A \( W^* \)-algebra is indeed a \textbf{locally convex topological vector space} when equipped with the weak-\( * \) topology (or even the norm topology). The weak-\( * \) topology, in particular, is locally convex and plays a central role in the theory of von Neumann algebras.

\section*{Question 5: What are the extreme points of the convex set of positive elements of a \( W^* \)-algebra?}

The set of \textbf{positive elements} in a \( W^* \)-algebra (von Neumann algebra) forms a convex cone, and understanding its \textbf{extreme points} is an important question in the study of operator algebras. Let's analyze this in detail:

\subsection*{1. Positive Elements in a \( W^* \)-Algebra:}
\begin{itemize}
    \item Let \( \mathcal{A} \) be a \( W^* \)-algebra. The set of positive elements in \( \mathcal{A} \), denoted \( \mathcal{A}_+ \), consists of all self-adjoint elements \( x \in \mathcal{A} \) with non-negative spectrum (\( \sigma(x) \subseteq [0, \infty) \)).
    \item \( \mathcal{A}_+ \) is a \textbf{convex cone}: it is closed under convex combinations and non-negative scalar multiplication.
\end{itemize}

\subsection*{2. Extreme Points of \( \mathcal{A}_+ \):}
\begin{itemize}
    \item An \textbf{extreme point} of a convex set \( C \) is a point \( x \in C \) that cannot be written as a non-trivial convex combination of two distinct points in \( C \). That is, if \( x = \lambda y + (1 - \lambda) z \) for \( y, z \in C \) and \( \lambda \in (0, 1) \), then \( y = z = x \).
    \item For the convex cone \( \mathcal{A}_+ \), the extreme points are the elements that cannot be expressed as a non-trivial convex combination of other positive elements.
\end{itemize}

\subsection*{3. Characterization of Extreme Points:}
\begin{itemize}
    \item The extreme points of \( \mathcal{A}_+ \) are precisely the \textbf{projections} in \( \mathcal{A} \). A projection \( p \in \mathcal{A} \) is an element satisfying \( p = p^* = p^2 \).
    \item \textbf{Why projections?}:
    \begin{itemize}
        \item Projections are positive (\( p \geq 0 \)) and idempotent (\( p^2 = p \)).
        \item If \( p \) is a projection and \( p = \lambda x + (1 - \lambda) y \) for \( x, y \in \mathcal{A}_+ \) and \( \lambda \in (0, 1) \), then \( x \) and \( y \) must also be projections, and \( x = y = p \). This shows that projections are extreme points.
        \item Conversely, if \( x \in \mathcal{A}_+ \) is not a projection, it can be decomposed into a non-trivial convex combination of other positive elements, so it is not an extreme point.
    \end{itemize}
\end{itemize}

\subsection*{4. Special Case: Finite-Dimensional \( W^* \)-Algebras:}
\begin{itemize}
    \item In the case where \( \mathcal{A} \) is a finite-dimensional \( W^* \)-algebra (i.e., a direct sum of matrix algebras \( M_n(\mathbb{C}) \)), the positive elements are positive semidefinite matrices.
    \item The extreme points of \( \mathcal{A}_+ \) are the \textbf{rank-one projections} (i.e., matrices of the form \( vv^* \), where \( v \) is a unit vector).
\end{itemize}

\subsection*{5. General Case: Infinite-Dimensional \( W^* \)-Algebras:}
\begin{itemize}
    \item In infinite-dimensional \( W^* \)-algebras, the extreme points of \( \mathcal{A}_+ \) are still the projections. However, the structure of projections can be more complex:
    \begin{itemize}
        \item Projections can be of finite or infinite rank.
        \item In a factor (a \( W^* \)-algebra with trivial center), projections are classified by their trace (in the case of a finite factor) or by their Murray-von Neumann equivalence class (in the case of an infinite factor).
    \end{itemize}
\end{itemize}

\subsection*{6. Geometric Interpretation:}
\begin{itemize}
    \item The set \( \mathcal{A}_+ \) is a convex cone, and its extreme points (the projections) lie on the "boundary" of the cone.
    \item The projections form a lattice under the partial order \( p \leq q \) if \( pq = p \), and this lattice structure is\section*{5. General Case: Infinite-Dimensional \( W^* \)-Algebras:}
\begin{itemize}
    \item In infinite-dimensional \( W^* \)-algebras, the extreme points of \( \mathcal{A}_+ \) are still the projections. However, the structure of projections can be more complex:
    \begin{itemize}
        \item Projections can be of finite or infinite rank.
        \item In a factor (a \( W^* \)-algebra with trivial center), projections are classified by their Murray-von Neumann equivalence class (in the case of an infinite factor).
    \end{itemize}
\end{itemize}

\section*{6. Geometric Interpretation:}
\begin{itemize}
    \item The set \( \mathcal{A}_+ \) is a convex cone, and its extreme points (the projections) lie on the "boundary" of the cone.
    \item The projections form a lattice under the partial order \( p \leq q \) if \( pq = p \), and this lattice structure is crucial in the study of \( W^* \)-algebras.
\end{itemize}

\section*{Conclusion:}
The extreme points of the convex set of positive elements \( \mathcal{A}_+ \) in a \( W^* \)-algebra are precisely the projections in \( \mathcal{A} \). These projections play a fundamental role in the structure and classification of \( W^* \)-algebras, and their study is closely tied to the spectral theory and order structure of operator algebras.

\section*{Question: Is the supremum of a continuous linear functional defined on the positive elements of a \( W^* \)-algebra only attained at projections?}

The question of whether the \textbf{supremum (sup)} of a continuous linear functional on the set of positive elements in a \( W^* \)-algebra (von Neumann algebra) occurs \textbf{only at projections} depends on the specific functional and the structure of the algebra. Let's analyze this carefully:

\subsection*{1. Setup:}
\begin{itemize}
    \item Let \( \mathcal{A} \) be a \( W^* \)-algebra, and let \( \mathcal{A}_+ \) be the set of positive elements in \( \mathcal{A} \).
    \item Let \( f: \mathcal{A} \to \mathbb{C} \) be a continuous linear functional. Since \( f \) is linear, it can be represented as \( f(x) = \phi(x) \) for some \( \phi \in \mathcal{A}_* \) (the predual of \( \mathcal{A} \)).
\end{itemize}

\subsection*{2. Supremum of \( f \) on \( \mathcal{A}_+ \):}
\begin{itemize}
    \item The supremum of \( f \) on \( \mathcal{A}_+ \) is defined as:
    \[
    \sup_{x \in \mathcal{A}_+} f(x).
    \]
    \item If \( f \) is continuous and \( \mathcal{A}_+ \) is weak-\( * \) compact (which is true if \( \mathcal{A} \) is unital), then the supremum is achieved at some point \( x \in \mathcal{A}_+ \).
\end{itemize}

\subsection*{3. Does the Supremum Occur Only at Projections?}
\begin{itemize}
    \item \textbf{Not necessarily}. The supremum of \( f \) on \( \mathcal{A}_+ \) can occur at elements that are \textbf{not projections}, depending on the functional \( f \).
    \item However, if \( f \) is a \textbf{normal positive linear functional} (i.e., \( f \in \mathcal{A}_* \) and \( f(x) \geq 0 \) for all \( x \in \mathcal{A}_+ \)), then the supremum of \( f \) on \( \mathcal{A}_+ \) is achieved at a projection if and only if \( f \) is \textbf{atomic} (i.e., it is associated with a minimal projection).
\end{itemize}

\subsection*{4. When Does the Supremum Occur at a Projection?}
\begin{itemize}
    \item If \( f \) is a \textbf{normal state} (a positive linear functional with \( f(1) = 1 \)), then the supremum of \( f \) on \( \mathcal{A}_+ \) is achieved at a projection if and only if \( f \) is \textbf{pure} (i.e., it cannot be written as a non-trivial convex combination of other states).
    \item In general, the supremum of \( f \) on \( \mathcal{A}_+ \) occurs at a projection \textbf{if and only if} \( f \) is \textbf{supported} on a projection. This means there exists a projection \( p \in \mathcal{A} \) such that \( f(x) = f(pxp) \) for all \( x \in \mathcal{A} \).
\end{itemize}

\subsection*{5. Counterexample:}
\begin{itemize}
    \item Consider the \( W^* \)-algebra \( \mathcal{A} = L^\infty([0, 1]) \), and let \( f \) be the linear functional defined by integration:
    \[
    f(x) = \int_0^1 x(t) \, dt.
    \]
    \item The supremum of \( f \) on \( \mathcal{A}_+ \) is \( 1 \), and it is achieved at the constant function \( x(t) = 1 \), which is \textbf{not a projection} in \( L^\infty([0, 1]) \).
\end{itemize}

\subsection*{6. Special Case: Finite-Dimensional \( W^* \)-Algebras:}
\begin{itemize}
    \item If \( \mathcal{A} \) is finite-dimensional (e.g., a matrix algebra \( M_n(\mathbb{C}) \)), then the supremum of \( f \) on \( \mathcal{A}_+ \) is achieved at a projection if and only if \( f \) is \textbf{rank-one} (i.e., it is associated with a pure state).
\end{itemize}

\subsection*{7. Conclusion:}
\begin{itemize}
    \item The supremum of a continuous linear functional \( f \) on \( \mathcal{A}_+ \) does \textbf{not necessarily occur only at projections}. It depends on the specific functional \( f \) and the structure of the algebra \( \mathcal{A} \).
    \item However, if \( f \) is a \textbf{normal positive linear functional} and the supremum is achieved at a projection, then \( f \) must be supported on that projection.
\end{itemize}

In summary, the supremum of a continuous linear functional on \( \mathcal{A}_+ \) can occur at elements that are not projections, but it occurs at a projection if and only if the functional is supported on that projection.